\documentclass[10pt,a4paper]{amsart}
\usepackage{amsmath,amssymb,amsthm,graphicx}
% Retained minimal preamble for the Stacks Project chapter varieties.tex.
% Verbatim excerpt of the chapter preamble: Xy-pic setup (native preamble lines 14-18), the native theorem
% declarations with the mathematical macros (native preamble lines 154-217), and the project own no-op
% definitions of the reference/slogan/history display environments (native wrapper.tex lines 1-3) that the
% counted statement uses for its slogan. Copied mechanically; nothing added, removed or redefined.
% For commutative diagrams we use Xy-pic
\usepackage[all]{xy}

% We use 2cell for 2-commutative diagrams.
\xyoption{2cell}
% Theorem environments.
%
\theoremstyle{plain}
\newtheorem{theorem}[subsection]{Theorem}
\newtheorem{proposition}[subsection]{Proposition}
\newtheorem{lemma}[subsection]{Lemma}

\theoremstyle{definition}
\newtheorem{definition}[subsection]{Definition}
\newtheorem{example}[subsection]{Example}
\newtheorem{exercise}[subsection]{Exercise}
\newtheorem{situation}[subsection]{Situation}

\theoremstyle{remark}
\newtheorem{remark}[subsection]{Remark}
\newtheorem{remarks}[subsection]{Remarks}

\numberwithin{equation}{subsection}

% Macros
%
\def\lim{\mathop{\mathrm{lim}}\nolimits}
\def\colim{\mathop{\mathrm{colim}}\nolimits}
\def\Spec{\mathop{\mathrm{Spec}}}
\def\Hom{\mathop{\mathrm{Hom}}\nolimits}
\def\Ext{\mathop{\mathrm{Ext}}\nolimits}
\def\SheafHom{\mathop{\mathcal{H}\!\mathit{om}}\nolimits}
\def\SheafExt{\mathop{\mathcal{E}\!\mathit{xt}}\nolimits}
\def\Sch{\mathit{Sch}}
\def\Mor{\mathop{\mathrm{Mor}}\nolimits}
\def\Ob{\mathop{\mathrm{Ob}}\nolimits}
\def\Sh{\mathop{\mathit{Sh}}\nolimits}
\def\NL{\mathop{N\!L}\nolimits}
\def\CH{\mathop{\mathrm{CH}}\nolimits}
\def\proetale{{pro\text{-}\acute{e}tale}}
\def\etale{{\acute{e}tale}}
\def\QCoh{\mathit{QCoh}}
\def\Ker{\mathop{\mathrm{Ker}}}
\def\Im{\mathop{\mathrm{Im}}}
\def\Coker{\mathop{\mathrm{Coker}}}
\def\Coim{\mathop{\mathrm{Coim}}}

% Boxtimes
%
\DeclareMathSymbol{\boxtimes}{\mathbin}{AMSa}{"02}

%
% Macros for moduli stacks/spaces
%
\def\QCohstack{\mathcal{QC}\!\mathit{oh}}
\def\Cohstack{\mathcal{C}\!\mathit{oh}}
\def\Spacesstack{\mathcal{S}\!\mathit{paces}}
\def\Quotfunctor{\mathrm{Quot}}
\def\Hilbfunctor{\mathrm{Hilb}}
\def\Curvesstack{\mathcal{C}\!\mathit{urves}}
\def\Polarizedstack{\mathcal{P}\!\mathit{olarized}}
\def\Complexesstack{\mathcal{C}\!\mathit{omplexes}}
% \Pic is the operator that assigns to X its picard group, usage \Pic(X)
% \Picardstack_{X/B} denotes the Picard stack of X over B
% \Picardfunctor_{X/B} denotes the Picard functor of X over B
\def\Pic{\mathop{\mathrm{Pic}}\nolimits}
\def\Picardstack{\mathcal{P}\!\mathit{ic}}
\def\Picardfunctor{\mathrm{Pic}}
\def\Deformationcategory{\mathcal{D}\!\mathit{ef}}
\newenvironment{reference}{}{}
\newenvironment{slogan}{}{}
\newenvironment{history}{}{}
\usepackage{cleveref}
\crefname{originalreference}{original reference}{original references}
\Crefname{originalreference}{Original reference}{Original references}
\makeatletter
\newlabel{situation-family-divisors}{{0G47}{source}{}{sourcepointer-1}{}}
\newlabel{situation-family-divisors@cref}{{[originalreference][0][]0G47}{[1][0][]source}{}{}{}}
\makeatother
\crefname{originallemma}{original lemma}{original lemmas}
\Crefname{originallemma}{Original lemma}{Original lemmas}
\makeatletter
\newlabel{lemma-dimension-fibres-locally-algebraic}{{0B2L}{source}{}{sourcepointer-2}{}}
\newlabel{lemma-dimension-fibres-locally-algebraic@cref}{{[originallemma][0][]0B2L}{[1][0][]source}{}{}{}}
\makeatother
\crefname{originalreference}{original reference}{original references}
\Crefname{originalreference}{Original reference}{Original references}
\makeatletter
\newlabel{morphisms-lemma-dimension-fibre-after-base-change}{{02FY}{source}{}{sourcepointer-3}{}}
\newlabel{morphisms-lemma-dimension-fibre-after-base-change@cref}{{[originalreference][0][]02FY}{[1][0][]source}{}{}{}}
\makeatother
\crefname{originalreference}{original reference}{original references}
\Crefname{originalreference}{Original reference}{Original references}
\makeatletter
\newlabel{morphisms-lemma-smooth-omega-finite-locally-free}{{02G1}{source}{}{sourcepointer-4}{}}
\newlabel{morphisms-lemma-smooth-omega-finite-locally-free@cref}{{[originalreference][0][]02G1}{[1][0][]source}{}{}{}}
\makeatother
\crefname{originalreference}{original reference}{original references}
\Crefname{originalreference}{Original reference}{Original references}
\makeatletter
\newlabel{algebra-lemma-rank-omega}{{00TR}{source}{}{sourcepointer-5}{}}
\newlabel{algebra-lemma-rank-omega@cref}{{[originalreference][0][]00TR}{[1][0][]source}{}{}{}}
\makeatother
\crefname{originalreference}{original reference}{original references}
\Crefname{originalreference}{Original reference}{Original references}
\makeatletter
\newlabel{divisors-lemma-fibre-Cartier}{{062Y}{source}{}{sourcepointer-6}{}}
\newlabel{divisors-lemma-fibre-Cartier@cref}{{[originalreference][0][]062Y}{[1][0][]source}{}{}{}}
\makeatother
\crefname{originalreference}{original reference}{original references}
\Crefname{originalreference}{Original reference}{Original references}
\makeatletter
\newlabel{morphisms-lemma-smooth-at-point}{{01V9}{source}{}{sourcepointer-7}{}}
\newlabel{morphisms-lemma-smooth-at-point@cref}{{[originalreference][0][]01V9}{[1][0][]source}{}{}{}}
\makeatother
\crefname{originalreference}{original reference}{original references}
\Crefname{originalreference}{Original reference}{Original references}
\makeatletter
\newlabel{algebra-lemma-regular-ring-CM}{{00NQ}{source}{}{sourcepointer-8}{}}
\newlabel{algebra-lemma-regular-ring-CM@cref}{{[originalreference][0][]00NQ}{[1][0][]source}{}{}{}}
\makeatother
\crefname{originalreference}{original reference}{original references}
\Crefname{originalreference}{Original reference}{Original references}
\makeatletter
\newlabel{algebra-lemma-separable-smooth}{{00TV}{source}{}{sourcepointer-9}{}}
\newlabel{algebra-lemma-separable-smooth@cref}{{[originalreference][0][]00TV}{[1][0][]source}{}{}{}}
\makeatother
\title{Bertini: the general member of a family of divisors cut by a generating system on a smooth projective immersion is smooth}
\author{The Stacks Project Authors}
\date{stacks-project:varieties.tex:local-snapshot-20261011; GNU FDL 1.2 or later}
\begin{document}
\maketitle
\noindent\textit{Editorial note.} All delivered mathematical text is the author's own, copied verbatim from the pinned Stacks Project chapter file varieties.tex; no statement or proof step was rewritten, reordered or removed.\par\medskip
\noindent\textit{Editorial note.} The proof contains commutative diagrams typeset with Xy-pic; the retained chapter preamble supplies the native Xy-pic setup, so they compile exactly as in the source.\par\medskip
\noindent\textit{Editorial note.} The counted span cites ten cross-chapter labels, each recorded in assembly.external\_references with its canonical tag as a navigation pointer.\par\medskip
\noindent\textit{Editorial note.} The proof cites no bibliography item inside the counted span.\par\medskip
\noindent\textit{Allowed standard prerequisites (editorial).}\begin{itemize}
\item The smooth locus of a morphism is open and fibres of smooth morphisms are smooth. Stacks Project, Morphisms, Lemma morphisms-lemma-smooth-locus-open (tag 01V9) and the definition of a smooth morphism. Makes the singular locus closed and reduces the claim to a statement about its image.
\item Dimension bounds for the singular locus of a family of divisors. Stacks Project, Varieties, section on Bertini theorems (the auxiliary dimension results used inside the proof). Supplies the estimate dim(Sigma) < r that finishes the argument.
\end{itemize}\medskip
\ifdefined\hypertarget\hypertarget{sourcepointer-1}{}\fi
\noindent\textit{Source reference.} Original reference 0G47: sources/stacks-chow/tags.txt, source lines 20117--20117. Cross-chapter standard prerequisite cited by the author inside the counted span; it is used as a pointer and its own proof is not part of this item.\par\medskip
\ifdefined\hypertarget\hypertarget{sourcepointer-2}{}\fi
\noindent\textit{Source reference.} Original lemma 0B2L: sources/stacks-chow/tags.txt, source lines 13934--13934. Cross-chapter standard prerequisite cited by the author inside the counted span; it is used as a pointer and its own proof is not part of this item.\par\medskip
\ifdefined\hypertarget\hypertarget{sourcepointer-3}{}\fi
\noindent\textit{Source reference.} Original reference 02FY: sources/stacks-chow/tags.txt, source lines 3209--3209. Cross-chapter standard prerequisite cited by the author inside the counted span; it is used as a pointer and its own proof is not part of this item.\par\medskip
\ifdefined\hypertarget\hypertarget{sourcepointer-4}{}\fi
\noindent\textit{Source reference.} Original reference 02G1: sources/stacks-chow/tags.txt, source lines 3212--3212. Cross-chapter standard prerequisite cited by the author inside the counted span; it is used as a pointer and its own proof is not part of this item.\par\medskip
\ifdefined\hypertarget\hypertarget{sourcepointer-5}{}\fi
\noindent\textit{Source reference.} Original reference 00TR: sources/stacks-chow/tags.txt, source lines 1085--1085. Cross-chapter standard prerequisite cited by the author inside the counted span; it is used as a pointer and its own proof is not part of this item.\par\medskip
\ifdefined\hypertarget\hypertarget{sourcepointer-6}{}\fi
\noindent\textit{Source reference.} Original reference 062Y: sources/stacks-chow/tags.txt, source lines 7812--7812. Cross-chapter standard prerequisite cited by the author inside the counted span; it is used as a pointer and its own proof is not part of this item.\par\medskip
\ifdefined\hypertarget\hypertarget{sourcepointer-7}{}\fi
\noindent\textit{Source reference.} Original reference 01V9: sources/stacks-chow/tags.txt, source lines 2438--2438. Cross-chapter standard prerequisite cited by the author inside the counted span; it is used as a pointer and its own proof is not part of this item.\par\medskip
\ifdefined\hypertarget\hypertarget{sourcepointer-8}{}\fi
\noindent\textit{Source reference.} Original reference 00NQ: sources/stacks-chow/tags.txt, source lines 868--868. Cross-chapter standard prerequisite cited by the author inside the counted span; it is used as a pointer and its own proof is not part of this item.\par\medskip
\ifdefined\hypertarget\hypertarget{sourcepointer-9}{}\fi
\noindent\textit{Source reference.} Original reference 00TV: sources/stacks-chow/tags.txt, source lines 1089--1089. Cross-chapter standard prerequisite cited by the author inside the counted span; it is used as a pointer and its own proof is not part of this item.\par\medskip
\begin{lemma}
\label{lemma-bertini}
In Situation \ref{situation-family-divisors} assume
\begin{enumerate}
\item $X$ is smooth over $k$,
\item the image of $\psi : V \to \Gamma(X, \mathcal{L})$
generates $\mathcal{L}$,
\item the corresponding morphism
$\varphi_{\mathcal{L}, \psi} : X \to \mathbf{P}(V)$ is
an immersion.
\end{enumerate}
Then for general $v \in V \otimes_k k'$
the scheme $H_v$ is smooth over $k'$.
\end{lemma}

\begin{proof}
(We observe that $X$ is separated and finite type as a locally
closed subscheme of a projective space.)
Let us use the notation introduced above the statement of the lemma.
We consider the projections
$$
\xymatrix{
\mathbf{A}^r_k \times_k X \ar[d] &
H_{univ} \ar[l] \ar[ld]^p \ar[r] \ar[rd]_q &
\mathbf{A}^r_k \times_k X \ar[d] \\
X & &
\mathbf{A}^r_k
}
$$
Let $\Sigma \subset H_{univ}$ be the singular locus of the morphsm
$q : H_{univ} \to \mathbf{A}^r_k$, i.e., the set of points where
$q$ is not smooth. Then $\Sigma$ is closed because the smooth locus
of a morphism is open by definition. Since the fibre of a
smooth morphism is smooth, it suffices to prove $q(\Sigma)$
is contained in a proper closed subset of $\mathbf{A}^r_k$.
Since $\Sigma$ (with reduced induced scheme structure) is a
finite type scheme over $k$ it suffices to prove $\dim(\Sigma) < r$
This follows from Lemma \ref{lemma-dimension-fibres-locally-algebraic}.
Since dimensions aren't changed by replacing $k$ by a bigger field
(Morphisms, Lemma \ref{morphisms-lemma-dimension-fibre-after-base-change}),
we may and do assume $k$ is algebraically closed.
By dimension theory (Lemma \ref{lemma-dimension-fibres-locally-algebraic}),
it suffices to prove that for $x \in X \setminus Z$ closed we have
$p^{-1}(\{x\}) \cap \Sigma$ has dimension $< r - \dim(X')$
where $X'$ is the unique irreducible component of $X$ containing $x$.
As $X$ is smooth over $k$ and $x$ is a closed point we have
$\dim(X') = \dim \mathfrak m_x/\mathfrak m_x^2$
(Morphisms, Lemma \ref{morphisms-lemma-smooth-omega-finite-locally-free}
and Algebra, Lemma \ref{algebra-lemma-rank-omega}).
Thus we win if
$$
\dim p^{-1}(x) \cap \Sigma < r - \dim \mathfrak m_x/\mathfrak m_x^2
$$
for all $x \in X$ closed.

\medskip\noindent
Since $V$ globally generated $\mathcal{L}$, for every irreducible component
$X'$ of $X$ there is a nonempty Zariski open of $\mathbf{A}^r$ such that the
fibres of $q$ over this open do not contain $X'$. (For example, if $x' \in X'$
is a closed point, then we can take the open corresponding to those
vectors $v \in V$ such that $\psi(v)$ does not vanish at $x'$. This
open will be the complement of a hyperplane in $\mathbf{A}^r_k$.)
Let $U \subset \mathbf{A}^r$ be the (nonempty) intersection of these opens.
Then the fibres of $q^{-1}(U) \to U$ are effective Cartier divisors
on the fibres of $U \times_k X \to U$ (because a nonvanishing section
of an invertible module on an integral scheme is a regular section). Hence the
morphism $q^{-1}(U) \to U$ is flat by
Divisors, Lemma \ref{divisors-lemma-fibre-Cartier}.
Thus for $x \in X$ closed and $v \in V = \mathbf{A}^r_k(k)$,
if $(x, v) \in H_{univ}$, i.e., if $x \in H_v$
then $q$ is smooth at $(x, v)$ if and only if the fibre
$H_v$ is smooth at $x$, see
Morphisms, Lemma \ref{morphisms-lemma-smooth-at-point}.

\medskip\noindent
Consider the image $\psi(v)_x$ in the stalk $\mathcal{L}_x$
of the section corresponding to $v \in V$. We have
$$
x \in H_v \Leftrightarrow \psi(v)_x \in \mathfrak m_x\mathcal{L}_x
$$
If this is true, then we have
$$
H_v\text{ singular at }x \Leftrightarrow
\psi(v)_x \in \mathfrak m_x^2\mathcal{L}_x
$$
Namely, $\psi(v)_x$ is not contained in $\mathfrak m_x^2\mathcal{L}_x$
$\Leftrightarrow$
the local equation for $H_v \subset X$ at $x$ is not contained
in $\mathfrak m_x^2$
$\Leftrightarrow$
$\mathcal{O}_{H_v, x}$ is regular
(Algebra, Lemma \ref{algebra-lemma-regular-ring-CM})
$\Leftrightarrow$
$H_v$ is smooth at $x$ over $k$
(Algebra, Lemma \ref{algebra-lemma-separable-smooth}).
We conclude that the closed points of $p^{-1}(x) \cap \Sigma$ correspond
to those $v \in V$ such that $\psi(v)_x \in \mathfrak m_x^2\mathcal{L}_x$.
However, as $\varphi_{\mathcal{L}, \psi}$ is an immersion the
map
$$
V \longrightarrow \mathcal{L}_x/\mathfrak m_x^2\mathcal{L}_x
$$
is surjective (small detail omitted). By the above, the closed points
of the locus $p^{-1}(x) \cap \Sigma$ viewed as a subspace of $V$
is the kernel of this map and hence has dimension
$r - \dim \mathfrak m_x/\mathfrak m_x^2 - 1$ as desired.
\end{proof}

\end{document}
